% !TeX root = ../main.tex

\section{Функциональное описание устройства}

Устройство выполняет аналого-цифровое преобразование зашумленного радиосигнала и его последующую цифровую квадратурную демодуляцию (перенос спектра почти на 0 частоту (из-за небольшого несоответсвия частот нельзя сказать что сигнал переносится на 0)).

\begin{itemize}
    \item \textbf{АЦП (Аналого-цифровой преобразователь):} Осуществляет дискретизацию входного аналогового сигнала с частотой $f_s$ и его квантование с разрядностью $n_a$ бит. В коде это реализовано функцией $quantize$ с жестким ограничением (насыщением) и округлением в меньшую сторону (Floor).
    \item \textbf{ЦГ (Цифровой гетеродин / NCO):} Формирует опорные квадратурные сигналы (косинусную и синусную составляющие) на частоте демодуляции $f_g$ с заданным разрешением $n_g$ бит. Модуль позволяет программно имитировать аппаратные неидеальности (рассогласование) приемника, вводя фазовый (\texttt{pmv}) и амплитудный (\texttt{amv}) дисбаланс между каналами.
    \item \textbf{Блок демодуляции (Перемножитель):} Осуществляет прямое умножение оцифрованного входного сигнала на комплексный сигнал цифрового гетеродина. На выходе формируются синфазная (I) и квадратурная (Q) составляющие сигнала, спектр которого сдвинут на величину частоты гетеродина.
\end{itemize}

В коде эти функции выполняют следующие функции:
\begin{itemize}
    \item \textbf{\texttt{get\_adc}}: Находится в конце скрипта. Имитирует работу аналого-цифрового преобразователя (АЦП). Выполняет квантование входного сигнала с заданной разрядностью \texttt{bits}, применяя округление в меньшую сторону (\texttt{Floor}) и ограничение сигнала при переполнении (\texttt{Saturate}).
    
    \item \textbf{\texttt{get\_nco}}: Формирует комплексный сигнал цифрового гетеродина и рассчитывает его коэффициент усиления (\texttt{nco\_gain}). Генерирует косинусную и синусную квадратурные составляющие, вносит в них искусственные амплитудные (\texttt{amv}) и фазовые (\texttt{pmv}) искажения, а также выполняет приведение типа выходных данных в форматы \texttt{fixed}, \texttt{single} или \texttt{double} в зависимости от переданного параметра.
    
    \item \textbf{\texttt{get\_dem}}: Моделирует блок демодуляции путем перемножения оцифрованного сигнала \texttt{adc\_signal} и комплексного сигнала гетеродина \texttt{nco\_signal}. Перед умножением функция приводит сигналы к нужному типу данных (\texttt{fixed}, \texttt{single}, \texttt{double}) и фиксирует суммарную разрядность результата для целочисленной арифметики.
    
    \item \textbf{\texttt{get\_spectrum}}:  Вычисляет спектр заданного сигнала с использованием быстрого преобразования Фурье (\texttt{fft}). Формирует вектор частот от $-f_s/2$ до $f_s/2$, переносит нулевую частоту в центр с помощью \texttt{fftshift} и опционально накладывает весовое окно Ханна (\texttt{hann}) перед БПФ, если параметр \texttt{weighting} равен 1.
\end{itemize}

Так же код был усовершенствован тем, что данные функции были вынесены из исновного кода в отдельные файлы для удобства отладки и наглядности кода. А также для индивидуальных значений задания по номеру студента была введена переменная N, для обозначения его номера в журнале.